\chapter{Kesimpulan dan Saran}

\section{Kesimpulan}

Penelitian ini telah menganalisis jaringan keilmuan Islam pada sanad hadis dengan memanfaatkan pendekatan \textit{Knowledge Graph} dan \textit{Social Network Analysis} (SNA), serta menyajikan hasil analisis tersebut melalui \textit{dashboard} interaktif berbasis web bagi mahasiswa dan peneliti ilmu hadis sebagai pengguna sasaran. Penelitian dilakukan melalui serangkaian tahapan, mulai dari pra-pemrosesan data sanad, pembangunan \textit{Knowledge Graph}, analisis jaringan menggunakan metrik sentralitas, hingga penyajian hasil analisis melalui \textit{dashboard} berbasis web. Berdasarkan hasil dan pembahasan yang telah diuraikan pada bab sebelumnya, dapat ditarik beberapa kesimpulan.

\begin{enumerate}
	\item Pembangunan \textit{Knowledge Graph} jaringan sanad hadis berhasil dilakukan berdasarkan relasi guru--murid yang diperoleh dari dataset Sanadset 650K. Data mentah sebanyak 650.986 baris terlebih dahulu melalui tahap pra-pemrosesan, meliputi penghapusan baris tanpa sanad serta normalisasi penulisan nama perawi seperti penghapusan harakat, penyeragaman variasi huruf Arab, dan resolusi kata kekerabatan. Tahap normalisasi ini penting karena memastikan satu perawi yang sama tidak terpecah menjadi beberapa simpul akibat perbedaan penulisan. Dari 487.451 rantai sanad yang bersih, terbentuk graf berarah yang terdiri atas 177.047 simpul perawi dan 776.798 sisi relasi periwayatan. Dengan demikian, jaringan keilmuan Islam yang sebelumnya hanya berupa daftar nama dapat dimodelkan sebagai jaringan terstruktur yang utuh dan dapat dianalisis secara komputasional.

	\item Peran dan pengaruh masing-masing perawi berhasil diidentifikasi melalui analisis SNA dengan lima metrik sentralitas, yaitu \textit{in-degree}, \textit{out-degree}, \textit{eigenvector}, \textit{closeness}, dan \textit{betweenness centrality}. Setiap metrik mencerminkan peran yang berbeda, yaitu perawi sebagai sumber riwayat (\textit{in-degree}), penerima riwayat (\textit{out-degree}), tokoh berpengaruh (\textit{eigenvector}), tokoh yang paling terjangkau (\textit{closeness}), serta penghubung antar-kelompok perawi (\textit{betweenness}). Hasil analisis menunjukkan bahwa Abu Hurairah menempati peringkat tertinggi pada sebagian besar metrik, yaitu \textit{in-degree}, \textit{eigenvector}, \textit{closeness}, dan \textit{betweenness}, sehingga teridentifikasi sebagai perawi yang paling berpengaruh dan paling sentral dalam jaringan, sedangkan Sufyan menempati peringkat tertinggi pada \textit{out-degree centrality}. Hal ini menunjukkan bahwa metrik sentralitas mampu membedakan peran tiap perawi secara terukur sekaligus menonjolkan perawi-perawi yang paling penting dalam jaringan periwayatan.

	\item Hasil analisis jaringan sanad berhasil disajikan melalui \textit{dashboard} web interaktif yang dibangun menggunakan React dengan data yang tersimpan pada basis data Supabase, lalu diterapkan secara daring sehingga dapat diakses publik tanpa memerlukan instalasi perangkat lunak khusus. Fitur pada \textit{dashboard} dirancang berdasarkan kebutuhan mahasiswa dan peneliti ilmu hadis sebagai pengguna sasaran (Subbab~\ref{subsubsec:kebutuhan-pengguna} Bab III), dan disajikan melalui empat halaman, yaitu Graf Relasi Perawi, Statistik Perawi, Tabel Relasi, dan Tabel Sentralitas. Halaman Graf Relasi Perawi dilengkapi visualisasi graf interaktif yang dirancang secara mandiri di atas HTML5 Canvas dengan algoritma tata letak \textit{force-directed}, sehingga perawi paling berpengaruh ditampilkan lebih menonjol dan berada di pusat jaringan. Hasil pengujian \textit{black-box} terhadap 14 skenario menunjukkan seluruh fitur berfungsi sesuai rancangan dengan tingkat keberhasilan 100\%.
\end{enumerate}

Berdasarkan ketiga hal tersebut, seluruh tujuan penelitian telah tercapai dan seluruh rumusan masalah telah terjawab berdasarkan data penelitian. Penelitian ini tidak hanya berhasil memodelkan dan menganalisis jaringan sanad hadis secara kuantitatif, tetapi juga menyajikan hasilnya dalam bentuk visualisasi interaktif yang mudah diakses oleh mahasiswa dan peneliti ilmu hadis, sehingga melengkapi keterbatasan penelitian terdahulu yang luarannya cenderung statis dan sulit dieksplorasi secara luas.

\section{Saran}

Berdasarkan hasil dan keterbatasan penelitian ini, terdapat beberapa saran yang dapat menjadi arah pengembangan pada penelitian selanjutnya.

\begin{enumerate}
	\item Perhitungan metrik sentralitas berbasis jalur, yaitu \textit{closeness} dan \textit{betweenness centrality}, pada penelitian ini masih terbatas pada subset graf terkuat karena keterbatasan komputasi pada perangkat yang digunakan. Penelitian selanjutnya disarankan untuk menghitung kedua metrik tersebut pada graf penuh, misalnya dengan memanfaatkan algoritma aproksimasi, komputasi paralel, atau perangkat dengan sumber daya yang lebih besar, agar nilai sentralitas yang diperoleh lebih menyeluruh dan akurat untuk seluruh perawi dalam jaringan.

	\item Data yang digunakan dalam penelitian ini masih bersifat statis dan belum diperbarui secara otomatis. Penelitian selanjutnya dapat mengembangkan mekanisme pembaruan data secara berkala atau mengintegrasikan sumber data sanad lain, sehingga jaringan yang dianalisis dapat terus berkembang dan tidak terbatas pada satu dataset saja.

	\item Analisis pada penelitian ini difokuskan pada struktur jaringan periwayatan dan belum mengkaji derajat kesahihan sanad secara ilmu hadis. Penelitian selanjutnya dapat menambahkan dimensi penilaian kualitas perawi maupun status kesahihan sanad, sehingga hasil analisis jaringan dapat dikaitkan langsung dengan kajian keilmuan hadis secara lebih mendalam.

	\item Validasi terhadap kebenaran setiap nama perawi juga belum dilakukan pada penelitian ini karena jumlah data yang sangat besar. Akibatnya, masih terdapat entri yang sebenarnya bukan nama perawi asli, melainkan kata umum atau sebutan yang merujuk pada makna lain, seperti kata yang berarti ``seorang laki-laki''. Penelitian selanjutnya disarankan untuk melakukan validasi dan pembersihan nama perawi secara lebih teliti, misalnya dengan bantuan kamus nama perawi atau verifikasi oleh ahli hadis, agar setiap simpul dalam jaringan benar-benar merepresentasikan perawi yang valid dan hasil analisis sentralitas menjadi lebih akurat.

	\item Evaluasi kebermanfaatan dan kegunaan \textit{dashboard} dari sisi pengguna belum dilakukan pada penelitian ini. Penelitian selanjutnya disarankan melakukan pengujian dengan mahasiswa dan peneliti ilmu hadis sebagai responden, misalnya melalui pengujian berbasis tugas (\textit{task-based testing}) dan kuesioner \textit{System Usability Scale} (SUS), untuk mengukur kebermanfaatan \textit{dashboard} secara empiris.
\end{enumerate}
